\section{Turning-angle curvature and the mixed-sign problem}\label{sec:discussion}

The choice of discrete curvature matters. Two elementary examples explain
why the results above neither extend to turning angles without further
hypotheses nor give a converse for arbitrary signed Menger profiles.

\subsection{Turning angles need not have four extrema}

If curvature is defined simply as the exterior turning angle, a strictly
convex polygon can have only one local maximum and one local minimum.

\begin{example}\label{ex:turning-angle}
Take the quadrilateral with counterclockwise vertices
\[
 A=(0,0),\qquad B=(3,0),\qquad C=(4,1),\qquad D=(0,2).
\]
It is strictly convex. Its exterior turns, in the same order, are
\[
 \theta_A=\frac{\pi}{2},\qquad
 \theta_B=\frac{\pi}{4},\qquad
 \theta_C=\frac{\pi}{2}+\arctan\frac35,\qquad
 \theta_D=\frac{\pi}{2}+\arctan\frac14.
\]
Thus $\theta_B<\theta_A<\theta_D<\theta_C$. Around the cycle
$A,B,C,D$, only $B$ is a local minimum and only $C$ is a local maximum.
All four readings are distinct, so plateau conventions do not affect the
example.
\end{example}

The construction is just a quadrilateral with unequal side lengths; there
is no limiting or degenerate configuration. It disproves the unrestricted
turning-angle version of the forward theorem. Additional assumptions can
restore such a statement: generic equilateral convex polygons with at least
four vertices satisfy an angle version of the discrete four-vertex theorem
\cite[Corollary~21.6]{Pak2010}.

\subsection{Four extrema do not suffice for mixed-sign Menger curvature}

For signed Menger curvature, the analogous difficulty concerns the
\emph{converse}. The strict four-vertex window alone does not guarantee
realizability once negative values are allowed.

\begin{example}\label{ex:mixed-sign}
For any $M,\varepsilon>0$, consider the cyclic profile
\begin{equation}\label{eq:mixed-counterexample}
 \kappa=(M,-\varepsilon,M,-\varepsilon).
\end{equation}
It has two strict maxima and two strict minima, and satisfies
$\DFV(\kappa,0)$: the two negative readings lie below zero and the two
positive readings above it. Nevertheless, no oriented simple quadrilateral
has this signed Menger profile.
\end{example}

\begin{proof}
For any counterclockwise simple polygon, let $\theta_i\in(-\pi,\pi)$ be its
principal exterior turns. They sum to $2\pi$. By~\eqref{eq:menger},
\[
 \kappa_i=\frac{2\sin\theta_i}{|P_{i+1}-P_{i-1}|},
\]
so $\kappa_i$ and $\theta_i$ have the same sign. There must therefore be at
least three positive curvature readings: with at most two positive turns,
\[
 2\pi=\sum_i\theta_i
 \le \sum_{\theta_i>0}\theta_i<2\pi,
\]
a contradiction. For a clockwise simple polygon the same argument with
signs reversed requires at least three negative readings. The profile
in~\eqref{eq:mixed-counterexample} has only two of each sign, so neither
orientation is possible.
\end{proof}

This obstruction is independent of coherence or local regularity. Making
$M$ arbitrarily large or $\varepsilon$ arbitrarily small does not remove it:
each individual positive turn is still less than $\pi$. Its moderate-arc
form is proved in \nolinkurl{Gluck/Discrete/Budget.lean}, by
\texttt{three\_le\_card\_pos\_of\_realizes} and
\texttt{killerProfile\_not\_realizes}. These Lean statements concern the
\texttt{RealizesR2} development model; the general polygon argument above
and Example~\ref{ex:turning-angle} are elementary arguments supplied here,
not separately formalized.

What additional conditions on a mixed-sign cyclic profile, beyond a strict
four-vertex window, characterize its realization as the signed Menger
curvature of a counterclockwise simple, locally regular polygon in the
sense of Dahlberg~\cite{Dahlberg2008}? The positive-support condition just
proved is necessary, but we do not claim it is sufficient. Thus the direct
mixed-sign extension is false, while finding the right strengthened
converse remains an open problem left by this work. Theorem~\ref{thm:main}
settles the positive case; a mixed-sign characterization would be a further
discrete counterpart of Dahlberg's smooth converse.
